%=============================================================================!
% 12.4  A SYSTEM AND ITS BATH                                                 !
%=============================================================================!
\section{A system and its bath}
\label{sec:sm-canonical}

Most systems of interest are not isolated: a sample exchanges energy with
its container, and a few hundred atoms in a simulation stand for a small
part of a much larger piece of matter. This section finds the
probabilities of the states of a system kept in contact with a much
larger one, its bath, and with them the partition function from which
averages and fluctuations follow.

\subsection*{The Boltzmann factor}

\begin{figure}[htb]
\centering
\begin{tikzpicture}[line cap=round, line join=round]
   \draw[line width=0.8pt, fill=black!6] (0,0) rectangle (6.4,2.6);
   \draw[line width=0.8pt, fill=white] (2.4,0.7) rectangle (4.0,1.9);
   \node[align=center] at (3.2,1.3) {system\\[2pt]\small $E_s$};
   \node[align=center] at (1.1,1.3) {bath\\[2pt]\small $E - E_s$};
   \node at (5.3,1.3) {\small energy};
   \draw[-{Stealth[length=5pt]}, line width=0.8pt, accent]
      (4.05,1.55) -- (4.75,1.55);
   \draw[{Stealth[length=5pt]}-, line width=0.8pt, accent]
      (4.05,1.05) -- (4.75,1.05);
\end{tikzpicture}
\caption{A small system inside a large bath. The two together are
isolated, with energy $E$; the system holds $E_s$ and the bath the rest,
and energy passes between them.}
\label{fig:sm-bath}
\end{figure}

The system and its bath (\cref{fig:sm-bath}) together are isolated, with
the energy $E$, so the pair is microcanonical: all its microstates of
energy $E$ are equally likely. We fix one microstate of the system, with
energy $E_s$. The bath then holds $E - E_s$, in any of its
$\Omega_{\mathrm B}(E - E_s)$ microstates, so the probability of that one
microstate of the system is proportional to $\Omega_{\mathrm B}(E - E_s)$.
The system's energy is small against the bath's, so we expand about $E$
in a Taylor series. $\Omega_{\mathrm B}$ itself changes enormously with its
energy, roughly as $E^{3N/2}$ for an ideal gas (\cref{sec:sm-micro}), but
its logarithm changes slowly, so it is the logarithm we expand:
\begin{align*}
   \ln\Omega_{\mathrm B}(E - E_s) &= \ln\Omega_{\mathrm B}(E)
   - E_s\,\frac{\dd\ln\Omega_{\mathrm B}}{\dd E}
   + \tfrac12E_s^2\,\frac{\dd^2\ln\Omega_{\mathrm B}}{\dd E^2} - \dotsb\\
   &= \ln\Omega_{\mathrm B}(E) - \frac{E_s}{\kB T} + \dotsb ,
\end{align*}
the second line by \cref{eq:sm-temperature}, with $T$ the temperature of
the bath. The second derivative is the rate
at which $\dd\ln\Omega_{\mathrm B}/\dd E = 1/\kB T$ changes, $-(1/\kB
T^2)\,\dd T/\dd E$ by the chain rule, so the next term is
$-\frac12E_s^2(\dd T/\dd E)/\kB T^2$; divided by the term kept,
$-E_s/\kB T$, it gives the ratio $\frac12E_s(\dd T/\dd E)/T$, in which
$\dd T/\dd E$ is how far the bath's temperature rises per unit of energy
it gains. For a bath that is an ideal gas, $T = 2E/3N\kB$ gives $\dd
T/\dd E = 2/3N\kB = T/E$, and the ratio is $E_s/2E$, half the system's
share of the energy: negligible for a bath much larger than the system
(\cref{ex:sm-bath}), and the dropped term itself, $(E_s/\kB T)(E_s/2E)$,
is small when $E \gg E_s^2/\kB T$. The factor $\Omega_{\mathrm B}(E)$ is the same for
every microstate of the system and joins the constant of proportionality,
so, taking the exponential of both sides,
\begin{equation}
   \mathcal{P}(\text{microstate}) \propto \mathrm{e}^{-E_s/\kB T}
   = \mathrm{e}^{-\beta E_s} , \qquad \beta = \frac{1}{\kB T} :
   \label{eq:sm-boltzmann}
\end{equation}
the probability of a state of the system falls exponentially with its
energy, by the Boltzmann factor. The bath enters only through its
temperature. The ensemble of a system at fixed $N$, $V$ and $T$ with these
probabilities is the canonical ensemble.

\subsection*{The partition function}

The probabilities must add up to one, so the probability of microstate
$s$ is $\mathcal{P}_s = \mathrm{e}^{-\beta E_s}/Z$, with $Z$ the sum over
all the system's microstates, or for the continuous phase space of
classical atoms the integral over it, the partition function:
\begin{equation}
   Z = \sum_s\mathrm{e}^{-\beta E_s} , \qquad\text{or}\qquad
   Z = \int\mathrm{e}^{-\beta H(\Gamma)}\,\dd\Gamma .
   \label{eq:sm-partition}
\end{equation}
with $\dd\Gamma$ measured in the unit of phase-space volume of
\cref{sec:sm-micro}, so that $Z$ is a pure number; changing that unit
adds a constant to $S$ and $\kB T$ times a constant to $F$, which drop out
of every difference and derivative below.
Averages follow from it by differentiation with respect to $\beta$:
\begin{align*}
   -\frac{\dd\ln Z}{\dd\beta} &= -\frac1Z\,\frac{\dd Z}{\dd\beta}
   = \frac{1}{Z}\sum_sE_s\mathrm{e}^{-\beta E_s} = \ens{E} ,\\
   \frac{\dd^2\ln Z}{\dd\beta^2} &= -\frac{\dd\ens{E}}{\dd\beta}
   = \frac{1}{Z}\sum_sE_s^2\mathrm{e}^{-\beta E_s}
   - \Bigl(\frac1Z\sum_sE_s\mathrm{e}^{-\beta E_s}\Bigr)^2
   = \ens{E^2} - \ens{E}^2 ,
\end{align*}
the first by the chain rule, with $\dd Z/\dd\beta = -\sum_sE_s
\mathrm{e}^{-\beta E_s}$, and the second by the product rule
(\cref{sec:ne-drag}) applied to $\ens{E} = (1/Z)\sum_sE_s
\mathrm{e}^{-\beta E_s}$: the sum has the derivative $-\sum_sE_s^2
\mathrm{e}^{-\beta E_s}$, and $1/Z$, by the chain rule, the derivative
$-(1/Z^2)\,\dd Z/\dd\beta = (1/Z^2)\sum_sE_s\mathrm{e}^{-\beta E_s}$.
The energy of a system in contact with a bath fluctuates, and its
variance is the rate at which $-\ens{E}$ changes with $\beta$. The heat
capacity $C_V = \dd\ens{E}/\dd T$ is the energy taken up per degree, at
fixed volume. By the chain rule $C_V = (\dd\ens{E}/\dd\beta)(\dd\beta/\dd
T)$, and $\dd\beta/\dd T = -1/(\kB T^2)$, so
\begin{equation}
   \ens{E^2} - \ens{E}^2 = -\frac{\dd\ens{E}}{\dd\beta}
   = \kB T^2\,C_V .
   \label{eq:sm-fluctuation}
\end{equation}
The size of the energy's fluctuations measures how much energy the
system takes up when warmed. The free energy is defined by
\begin{equation*}
   F = -\kB T\ln Z ,
\end{equation*}
and equals $\ens{E} - TS$ with the entropy $S = -\kB\sum_s\mathcal{P}_s\ln
\mathcal{P}_s$, which for the equal probabilities $1/\Omega$ of the
microcanonical ensemble is \cref{eq:sm-entropy} (\cref{ex:sm-free}).

\subsection*{A spring, and the odds of a hop}

One coordinate on a spring, $H = p^2/2m + \frac12kx^2$, has, by the
Gaussian integral of \cref{sec:md-ewald} twice,
\begin{equation*}
   Z = \int\mathrm{e}^{-\beta p^2/2m}\dd p\int\mathrm{e}^{-\beta kx^2/2}
   \dd x = \sqrt{\frac{2\pi m}{\beta}}\sqrt{\frac{2\pi}{\beta k}}
   = \frac{2\pi}{\beta\omega} ,
\end{equation*}
with $\omega = \sqrt{k/m}$. Then $\ln Z = -\ln\beta + \text{a constant}$,
and $\ens{E} = -\dd\ln Z/\dd\beta = 1/\beta = \kB T$: a classical vibration holds
$\kB T$ on average, whatever its stiffness, which makes precise the
`typical energy of order $\kB T$' of \cref{sec:pe-classical}. At
\SI{300}{\kelvin}, $\kB T = \SI{0.02585}{\electronvolt}$.

The Boltzmann factor also answers the question of \cref{sec:en-surface},
how often an atom reaches the top of its barrier. The kinetic energy
depends on the momenta alone and $U$ on the positions alone, so
$\mathrm{e}^{-\beta H}$ is $\mathrm{e}^{-\beta U}$ times a factor of the
momenta, and integrating over the momenta turns that factor into a number
that is the same at every position. The positions are left with the
density $\mathrm{e}^{-\beta U}$, in which the masses no longer appear, as
\cref{sec:cs-mass} anticipated: of two places whose potential energies
differ by $\Delta U$, an atom is found near the higher, per unit volume,
$\mathrm{e}^{-\Delta U/\kB T}$ times as often as near the lower. For the
lithium atom of that model, which moves over the surface, the density is
per unit area, and at the bridge, \SI{0.3}{\electronvolt} above its
hollow, the odds are $\mathrm{e}^{-0.3/0.02585} = \num{9.1e-6}$ at
\SI{300}{\kelvin}, and $\num{3.0e-3}$ at \SI{600}{\kelvin}, 331 times
larger: a barrier that is rarely reached at room temperature is reached
hundreds of times as often when the temperature doubles. For an atom
moving in two directions the same number is also the probability that its
kinetic energy is at least the barrier, since its kinetic energy has the
density $\mathrm{e}^{-K/\kB T}/\kB T$ (\cref{sec:sm-kinetic}). How often an atom actually hops depends also on how
often it tries, which Chapters~16 and~17 measure.

\begin{keyidea}
A system in contact with a bath at temperature $T$ is found in a
microstate of energy $E_s$ with a probability proportional to
$\mathrm{e}^{-E_s/\kB T}$. The partition function $Z$ normalises the
probabilities; $\ens{E} = -\dd\ln Z/\dd\beta$, the energy's variance is
$\kB T^2C_V$, and $F = -\kB T\ln Z$. A classical vibration holds $\kB T$
on average.
\end{keyidea}

\notebook{12}{set a barrier and a temperature and read off the odds}

\begin{selfcheck}
What are the odds of the \SI{0.3}{\electronvolt} barrier at
\SI{100}{\kelvin}?
\end{selfcheck}
